\section{Appendix}
\subsection{Type Conversion: \texttt{dict()} and \texttt{list()}}
\begin{itemize}
    \item \textbf{Tuple to Dict}: The \texttt{dict()} function can convert a list of 2-item tuples directly into key-value pairs.
    \item \textbf{Deduplication}: Since dictionary keys must be unique, converting a list of tuples to a \texttt{dict} automatically removes duplicate entries (keeping the last evaluated value for a given key).
    \item \textbf{Dict to List}: Converts the deduplicated dictionary back into a list of tuples using \texttt{list(dictionary.items())}.
\end{itemize}

\lstinputlisting[language=Python]{code/type_coversion.py}

\subsection{Decorators}
Decorators are functions that take another function as an argument to extend its behavior without modifying the original source code. They are denoted by the \texttt{@} symbol.

\begin{lstlisting}[language=Python]
def require_login(func):
    def wrapper(user):
        if user["logged_in"]:
            return func(user)
        print("Access denied.")
    return wrapper

@require_login
def view_dashboard(user):
    print(f"Welcome {user['name']}")
\end{lstlisting}


\subsection{Script Execution}
To distinguish between a script being run directly and a script being imported as a module, Python uses the \texttt{\_\_name\_\_} variable.

\begin{itemize}
    \item When run directly: \texttt{\_\_name\_\_ == "\_\_main\_\_"}.
    \item When imported: \texttt{\_\_name\_\_} equals the module's filename.
\end{itemize}

\begin{lstlisting}[language=Python]
def main():
    # Test code or entry point
    pass

if __name__ == "__main__":
    main()
\end{lstlisting}

\subsection{Distribuzioni Casuali in NumPy}

\begin{itemize}
    \item \textbf{np.random.rand}: Genera campioni da una distribuzione \textbf{uniforme} nell'intervallo $[0, 1)$. Ogni valore ha la stessa probabilità di essere estratto.
    \item \textbf{np.random.randn}: Genera campioni dalla distribuzione \textbf{normal standard} (Gaussiana) con media $\mu = 0$ e varianza $\sigma^2 = 1$.
\end{itemize}

\subsubsection{Trasformazione della Distribuzione Normale}
Per scalare e traslare la distribuzione normale standard $\mathcal{N}(0, 1)$ verso una distribuzione personalizzata $\mathcal{N}(\mu, \sigma^2)$, si applica la seguente formula lineare:

\[ X = \sigma \cdot \text{np.random.randn}(\dots) + \mu \]

Dove:
\begin{itemize}
    \item $\sigma$ (Sigma): Deviazione standard (controlla la "larghezza" della campana).
    \item $\mu$ (Mu): Media (sposta il centro della distribuzione).
\end{itemize}

\subsection{Indexing vs Slicing: Regola Rapida}

La differenza dipende dal tipo di selettore utilizzato per l'asse:

\begin{itemize}
    \item \textbf{Intero (Indexing)}: Se usi un numero secco (es. \texttt{arr[2]}), la dimensione viene \textbf{rimossa} (rank-1).
    \item \textbf{Range (Slicing)}: Se usi i due punti (es. \texttt{arr[2:3]}), la dimensione viene \textbf{mantenuta} (rank-2).
\end{itemize}

\subsubsection{Esempio Pratico}
Dato un array 2D di shape \texttt{(3, 3)}:
\begin{itemize}
    \item \texttt{arr[2, :]} $\rightarrow$ Risultato 1D, shape \texttt{(3,)}.
    \item \texttt{arr[2:3, :]} $\rightarrow$ Risultato 2D, shape \texttt{(1, 3)}.
\end{itemize}